Plots kreeg ze een idee.

De jager was strafbaar
na wat hij deed met haar.
Melden was weliswaar

vooral voor haar een ramp -
immers het bijbelkamp
schoot altijd in de kramp

bij morele taken,
die hadden te maken
met geslachtsdriftzaken -

maar belastte ook hem
en zette hem ook klem.
Al had een vrouwenstem

zonder reputatie
niet veel relevantie,
bij een impregnatie

had ze zichtbaar bewijs.
Echter, het was niet wijs
om een denkbare eis

juridisch te halen.
Al zulke verhalen
lieten haar afdwalen

van het bloot rondlopen.
Ze begon te hopen
om het te ontlopen.

Eerst was ze slechts marktwaar.
Toen ontvoerde hij haar.
Ook dat was toen strafbaar.

Net als haar eerste hulp,
die nog met open gulp
net ontdaan van zijn pulp,

stoer voor Emmy opkwam
tot men zijn leven nam.
Met dit intern gezwam

waande ze zich veilig.
De jager had geilig
ook Emmy schijnheilig

geholpen tegen kwaad.
Hij had haar na zijn daad,
immers nog meer geschaad,

als hij haar bezwangerd
en reinheid geplunderd,
toch had uitgeleverd.

Hij had groots gezworen
dat zulks een ontsporen
volstrekt niet zou horen

bij zijn bedoelingen.
Bij zijn binnendringen
liet hij soms doordringen

in enkele zuchten
dat hij haar wou tuchten
om haar te bevruchten.

Wilde hij graag een kind?
Had geen vrouw welgezind
hem voldoende bemind

met hem te gaan paren
om zijn kind te baren?
Wilde hij haar sparen

door met haar te huwen
om dan zijn schaduwen
in haar buik te duwen?

Dat beloofde wel wat!
In een andere stad
had de landheer geen vat

op haar vermoedelijk.
Maakte een huwelijk
het Emmy mogelijk

anoniem te blijven.
In elk formeel schrijven
kwam niet bovendrijven

haar echte achternaam.
Dan had Emmy bekwaam
een aanwijzing tot blaam

weten uit te vegen.
Een naamloze zegen.
Ze bleef overwegen.

Hij werd haar echtgenoot.
Waar ze niet van genoot,
was zijn huwelijksboot.

Hij was immers lelijk,
verslaafd en ouwelijk.
Trouwen leek gruwelijk,

maar uit haar opvoeding
kwam de overweging
dat alleen hun paring

door echtverbintenis
kwam tot vergiffenis.
Ondanks de ergernis

over de nadelen
van het bed te delen,
kon ze toch meespelen.

Werd niet de nymph Chloris
uit de geschiedenis
van de Griekse mythes

door Zephyrus ontvoerd,
met de wind weggevoerd
en toen onkuis beroerd.

Er ontstond harmonie
in deze relatie,
ondanks de infractie.

De levensgezellin
werd beloond als godin.
Emmy zag er brood in

juist toen ze het koud kreeg,
nadat hij haar besteeg
en lijven aaneenreeg.

Hij zou niet verzaken
over haar te waken
na haar zwanger raken,

hoewel dat natuurlijk
nog alleen wenselijk
was en niet werkelijk.

Dus na enige schroom
keek Emmy quasi vroom
om naar de denneboom

haar bedekking missend.
De jager, al pissend
wat dauwdruppels sissend,

verwierp Emmy’s oproep. “Omkeren voor die troep?
Ik wil een bordje soep,

dus ik wil naar de stad.
Dus beweeg snel je gat,
want anders zwaait er wat!

Jij ondankbare trut.
Beweeg je gratenkut.”
Zo schold hij haar voor schut.

Emmy dacht opgeschrikt:
"Alles heb ik gepikt.
Spuug, zweet en zaad geslikt.

Je scheten geroken.
Je vuur opgestoken.
Mijn vlies is gebroken.

Mijn buikje vederlicht
was nagenoeg gezwicht
voor jouw over(ge)wicht.

Je mocht me uitschelden.
Mijn kastijding melden.
Ik moest het ontgelden.

Alles om jouw climax,
die je, kreunend daarstraks
barbaars en achterbaks,

in een tienermeisje
inpaste, als prijsje
om diens levensreisje

drastisch te verdraaien
door een kind te zaaien.
Jou levenslang naaien

als echtelijke plicht
is nu het vooruitzicht
voor mij als puberwicht.

Je hoort mij niet klagen
na dit onbehagen.
Ik wil je wel vragen,

me uitgekookt geschaakt,
me tot je vrouw gemaakt,
om me niet langer naakt

te laten rondlopen.
Ik mag toch wel hopen,
al ben je bezopen,

dat je niet zou willen
dat iemand mijn billen
ontkleed kan zien trillen?”

Deze overweging
als lange opsomming
gaf de overtuiging

om te vragen met moed
naar haar billijk tegoed.
Dat viel bij hem niet goed,

al koos ze haar woorden
zoals ze behoorden.
Zijn flaporen hoorden

echter alleen opstand.
Hij verbeet de weerstand
door met een platte hand

te meppen op haar kont.
Verstandelijk verwond,
hield ze meteen haar mond

na een krassende gil.

Hij wierp haar een schampere lach toe toen ze hem om begrip vroeg.
